% === Section 5: Contracts and Governance ===
\section{Contract Stack and Governance Surface}\label{sec:contracts}

Every \dcs{} step produces intermediate typed objects.  The \textbf{contract stack}
defines the obligations that each object must satisfy for the step to be valid.  The
\textbf{governance surface} aggregates per-step certificates into episode- and
campaign-level summaries that auditors can inspect.

\subsection{Typed Contract Objects}

\Cref{tab:contract_objects} lists the principal contract objects, their
associated assumptions, and their role in the kernel.

\begin{table}[t]
\caption{Typed contract objects produced by the \dcs{} kernel.}
\label{tab:contract_objects}
\centering\scriptsize
\setlength{\tabcolsep}{2.5pt}
\renewcommand{\arraystretch}{1.08}
\begin{tabularx}{\linewidth}{p{2.7cm}p{1.25cm}X}
\toprule
\textbf{Contract object} & \textbf{A} & \textbf{Role} \\
\midrule
\texttt{ObservationPacket}            & A1       & Raw telemetry + parsed state \\
\texttt{ReliabilityAssessment}        & A2, A5, A6 & Detection output: $\wt$, flags \\
\texttt{ObsConsistentStateSet} & A2, A5   & Uncertainty envelope $\Ut$ \\
\texttt{SafetySpec}                   & A3, A4, A8 & Domain constraints + fallback \\
\texttt{SafeActionSet}                & A3, A4   & Tightened admissible actions \\
\texttt{RepairDecision}              & A4       & Shield output \\
\texttt{SafetyCertificate}           & CP1--CP6 & Auditable record binding all stages \\
\bottomrule
\end{tabularx}
\end{table}

\subsection{ContractVerifier}

The \texttt{ContractVerifier} is a second-pass checker that validates a
\texttt{SafetyCertificate} against a set of \emph{contract predicates}.  Each
predicate is a boolean function of the certificate's payload.  The verifier evaluates
all predicates and returns a \texttt{VerificationResult} with:
\begin{itemize}[nosep]
  \item \textbf{Chain validity}: conjunction of all predicate outcomes.
  \item \textbf{Failure rows}: count of certificates that failed at least one predicate.
  \item \textbf{Required-payload rate}: fraction of certificates with all
        required fields present.
  \item \textbf{Audit completeness}: fraction of steps with a valid certificate.
\end{itemize}

\subsection{Contract Predicates}

The predicate set is domain-agnostic.  The following predicates are evaluated for
every certificate:

\begin{enumerate}[label=\textbf{CP\arabic*},nosep]
  \item \textbf{Monotone inflation}: the inflation factor $\alpha \ge 1.0$.
  \item \textbf{Weight bounds}: $\wt \in [0, 1]$.
  \item \textbf{Interval non-degenerate}: for every target, $\mathrm{upper} \ge
        \mathrm{lower}$.
  \item \textbf{Action within bounds}: released action satisfies all tightened
        constraints.
  \item \textbf{Certificate completeness}: all required fields are present and
        non-null.
  \item \textbf{Timestamp monotonicity}: certificate timestamps are non-decreasing
        within an episode.
\end{enumerate}

The contract predicate CP1 (monotone inflation) is a runtime-law check: lower
reliability may widen the uncertainty set, but it must never shrink it.  The
observation-ambiguity result in Theorem~T4 then explains why this check matters:
when degraded observations collapse distinct hidden states into the same
observation class, quality-ignorant release can become unsafe.

\subsection{Domain Instantiation Protocol}

A \texttt{DomainInstantiation} adapter implements the small typed surface in
\cref{tab:domain_protocol_methods}.  This keeps the kernel stable: adding a
domain means implementing these methods, not changing the release contract.

\begin{table}[t]
\caption{\texttt{DomainInstantiation} adapter surface.}
\label{tab:domain_protocol_methods}
\centering\footnotesize
\begin{tabularx}{\linewidth}{lX}
\toprule
\textbf{Method} & \textbf{Contract output} \\
\midrule
\texttt{domain\_id} & Stable domain identifier for evidence rows and certificates. \\
\texttt{assess} & \texttt{ReliabilityAssessment} from current observation and history. \\
\texttt{calibrate} & \texttt{ObservationConsistentStateSet} with reliability-aware uncertainty. \\
\texttt{constrain} & \texttt{SafeActionSet} induced by the uncertainty set and safety spec. \\
\texttt{shield} & \texttt{RepairDecision}: release, repair, fallback, or deny. \\
\texttt{certify} & \texttt{SafetyCertificate} binding the released action to evidence. \\
\bottomrule
\end{tabularx}
\end{table}

\subsection{Governance Surface}\label{sec:governance_surface}

At the episode level, the governance surface aggregates certificates into a summary
\texttt{runtime\_governance\_summary.csv} with the metrics in
\cref{tab:governance_metrics}.

\begin{table*}[t]
\caption{Claim-governing governance summary metrics regenerated from the active
certificate artifacts.}
\label{tab:governance_metrics}
\centering\small
\begin{tabular}{lccc}
\toprule
\textbf{Metric} & \textbf{Battery} & \textbf{nuPlan AV} & \textbf{Healthcare} \\
\midrule
\texttt{certificate\_rows} & 1{,}152 & 1{,}531{,}104 & 137{,}262 \\
\texttt{chain\_valid} & 1.0 & 1.0 & 1.0 \\
\texttt{failure\_rows} & 0 & 0 & 0 \\
\texttt{required\_payload\_pass\_rate} & 1.0 & 1.0 & 1.0 \\
\texttt{expiry\_metadata\_presence\_rate} & 1.0 & 1.0 & 1.0 \\
\texttt{expiry\_consistency\_rate} & 1.0 & 1.0 & 1.0 \\
\texttt{audit\_completeness\_rate} & 1.0 & 1.0 & 1.0 \\
\bottomrule
\end{tabular}
\end{table*}

The battery governance row covers all four promoted battery runtime variants
(1{,}152 certificates).  The all-zip grouped nuPlan AV row reports a 99.992\%
certificate-valid release rate and 100\% audit completeness on 1{,}531{,}104
runtime steps.

\subsection{CertOS Verification}

The CertOS module performs an offline re-verification of the entire certificate chain.
For each active runtime certificate surface:
\begin{itemize}[nosep]
  \item \textbf{Battery}: 1{,}152 certificates verified, \texttt{chain\_valid = true},
        0 failures.
  \item \textbf{AV}: 1{,}531{,}104 all-zip grouped nuPlan certificates emitted,
        certificate-valid release rate $0.999924$,
        249 empirical runtime violations under the promoted replay contract.
\end{itemize}

\subsection{Safety Architecture Comparison}

\Cref{tab:safety_comparison} positions ORIUS/\dcs{} against five alternative
safety-layer paradigms.  The comparison focuses on structural capabilities rather
than domain-specific numbers.

\begin{table*}[t]
\caption{Safety architecture comparison. Entries are structural support levels,
not trainable metrics: \emph{Yes} means native support, \emph{Partial} means
support under additional modeling assumptions, and \emph{No} means the baseline
does not natively expose that ORIUS release-contract obligation.}
\label{tab:safety_comparison}
\centering\scriptsize
\setlength{\tabcolsep}{3.2pt}
\renewcommand{\arraystretch}{1.08}
\begin{tabularx}{0.98\textwidth}{p{3.1cm}XXXXXX}
\toprule
\textbf{Capability} & \textbf{DC3S} & \textbf{Simplex} & \textbf{GP-MPC}
  & \textbf{DRO} & \textbf{Safe RL} & \textbf{Vanilla CQR} \\
\midrule
Obs-quality aware & Yes & No & No & No & No & No \\
Reliability-adaptive margin & Yes & No & Partial & No & No & No \\
Conformal coverage guarantee & Yes & No & No & No & No & Yes \\
Domain-agnostic kernel & Yes & Partial & No & No & No & Yes \\
Typed adapter protocol & Yes & No & No & No & No & No \\
Certificate governance & Yes & No & No & No & No & No \\
Graceful degradation (T8) & Yes & Yes & Partial & No & No & No \\
Formal proof programme & Yes & Yes & Partial & Yes & Partial & Partial \\
Runtime overhead $<$ 1\,ms & Yes & Yes & No & No & No & Yes \\
Three-row validated\textsuperscript{$\dagger$} & Yes & No & No & No & No & No \\
\bottomrule
\end{tabularx}
\end{table*}

{\footnotesize $\dagger$ Active defended lane: Battery witness row plus bounded AV and Healthcare rows with locked artifacts.}

\paragraph{Key differentiators.}
\begin{itemize}[nosep]
  \item \textbf{Simplex}~\cite{sha1998simplex,sha2001using} provides safety switching between a
        performance controller and a verified safe controller, but does not adapt
        the safety margin to telemetry reliability.
  \item \textbf{GP-MPC} (Gaussian Process Model Predictive Control) provides
        uncertainty-aware control but requires a Gaussian process model per domain,
        with $O(n^3)$ scaling precluding real-time use at large state dimensions.
  \item \textbf{DRO} (Distributionally Robust Optimisation) handles distributional
        ambiguity but does not produce per-step safety certificates or adapt to
        observation-quality signals.
  \item \textbf{Safe RL} incorporates safety constraints into the reward, but
        typically provides asymptotic guarantees rather than per-step certificates.
  \item \textbf{Vanilla CQR} achieves conformal coverage but is quality-ignorant:
        the same interval width is used regardless of telemetry degradation.
\end{itemize}

DC3S is unique in combining observation-quality awareness, reliability-adaptive margins,
a formal proof programme, domain-agnostic typed adapters, and per-step certificate
governance in a single runtime layer with sub-millisecond overhead.
